\documentclass[../../main.tex]{subfiles}
\begin{document}
\section{齐性空间}\label{section:ch21-homogeneous-spaces}

最有趣的群作用是可迁作用。赋予 Lie 群 $G$ 的可迁光滑作用的光滑流形称为\emph{齐性 $G$-空间}（若无需指明群，简称齐性空间或齐性流形）。

\begin{example}{[齐性空间]}\label{example:ch21-齐性空间}
(a) $O(n)$ 在 $\mathbb{S}^{n-1}$ 上的自然作用可迁。
(b) 把 $O(n)$ 的自然作用限制到 $\operatorname{SO}(n)$ 在 $\mathbb{S}^{n-1}$ 上的光滑作用。$n=1$ 时平凡；$n>1$ 时可迁：对任意 $v\in\mathbb{S}^{n-1}$，存在 $A\in O(n)$ 把第一标准基向量 $e_1$ 映到 $v$；若 $\det A=-1$，把 $A$ 第二列乘 $-1$ 即得 $\operatorname{SO}(n)$ 中矩阵。故 $n\geqslant2$ 时 $\mathbb{S}^{n-1}$ 也是 $\operatorname{SO}(n)$ 的齐性空间。
(c) 例 7.32 的 Euclid 群 $E(n)$ 通过刚体运动作用于 $\mathbb{R}^n$。平移可把任意点映到任意点，故 $E(n)$ 在 $\mathbb{R}^n$ 上可迁，$\mathbb{R}^n$ 是齐性 $E(n)$-空间。
(d) $\operatorname{SL}(2,\mathbb{R})$ 通过 $\begin{pmatrix}a&b\\ c&d\end{pmatrix}\cdot z=\dfrac{az+b}{cz+d}$ 光滑可迁作用于上半平面 $U=\{z\in\mathbb{C}\colon\operatorname{Im}z>0\}$。这些从 $U$ 到自身的复解析微分同胚称为 Möbius 变换。
(e) $n\geqslant1$ 时 $\operatorname{GL}(n,\mathbb{C})$ 在 $\mathbb{C}^n$ 上的自然作用限制为 $U(n)$ 与 $\operatorname{SU}(n)$ 在 $\mathbb{S}^{2n-1}$（视为 $\mathbb{C}^n$ 中单位向量集）上的自然光滑作用。
\end{example}

\begin{exercise}\label{exercise:ch21-Un与SUn可迁}
证明 $U(n)$ 在 $\mathbb{S}^{2n-1}$ 上的自然作用对所有 $n\geqslant1$ 可迁，$\operatorname{SU}(n)$ 的作用对所有 $n\geqslant2$ 可迁。
\end{exercise}

下面描述一个可用来由 Lie 群对闭子群的商生成大量齐性空间的构造。设 $G$ 是 Lie 群，$H\subseteq G$ 是 Lie 子群。形如 $gH=\{gh\colon h\in H\}$（$g\in G$）的子集称为 $H$ 的\emph{左陪集}。左陪集划分 $G$，由该划分确定的商空间（左陪集集合赋商拓扑）称为 $G$ 模 $H$ 的\emph{左陪集空间}，记为 $G/H$。$g_1,g_2\in G$ 在同一左陪集当且仅当 $g_1^{-1}g_2\in H$；此时记 $g_1\equiv g_2\pmod H$。

\begin{theorem}{[齐性空间构造定理]}\label{theorem:ch21-齐性空间构造}
设 $G$ 是 Lie 群，$H$ 是 $G$ 的闭子群。则左陪集空间 $G/H$ 是维数为 $\dim G-\dim H$ 的拓扑流形，并有唯一光滑结构使商映射 $\pi\colon G\to G/H$ 是光滑淹没。$G$ 在 $G/H$ 上由 $g_0\cdot(gH)=(g_0 g)H$ 给出的左作用光滑、可迁。
\end{theorem}

\begin{proof}
令 $H$ 通过右平移作用于 $G$。则 $g_1,g_2$ 在同一 $H$-轨道当且仅当 $g_1 h=g_2$（某 $h\in H$），即同一 $H$ 陪集。故 $H$ 在 $G$ 上右作用确定的轨道空间恰为左陪集空间 $G/H$。由闭子群定理，$H$ 是 $G$ 的真嵌入 Lie 子群；$H$ 在 $G$ 上的作用光滑（只是 $G$ 乘法的限制）。它自由：$gh=g$ 蕴含 $h=e$。为证真，用命题\cref{proposition:ch21-真作用刻画}(b)：若 $(g_i)$ 在 $G$ 中收敛、$(h_i)\subset H$ 使 $(g_i h_i)$ 收敛，则由连续性 $h_i=g_i^{-1}(g_i h_i)$ 收敛于 $G$ 中一点；因 $H$ 在 $G$ 中闭且带子空间拓扑，$(h_i)$ 在 $H$ 中收敛。
商流形定理于是给出 $G/H$ 唯一光滑流形结构使 $\pi$ 是光滑淹没。因光滑淹没的乘积是光滑淹没，$\operatorname{Id}_G\times\pi\colon G\times G\to G\times(G/H)$ 也是光滑淹没。考虑图表 $G\times G\xrightarrow{m}G$ 与 $G\times(G/H)\xrightarrow{\theta}G/H$，其中 $m$ 为群乘法、$\theta$ 为上述左作用。易验 $\theta\circ(\operatorname{Id}_G\times\pi)$ 在 $\operatorname{Id}_G\times\pi$ 的纤维上常值，故由定理 4.30，$\theta$ 良定义且光滑。由定义它满足群作用条件。最后，对任意 $g_1H,g_2H\in G/H$，元素 $g_2g_1^{-1}\in G$ 满足 $(g_2g_1^{-1})\cdot(g_1H)=g_2H$，故作用可迁。
\end{proof}

\begin{theorem}{[齐性空间刻画定理]}\label{theorem:ch21-齐性空间刻画}
设 $M$ 是光滑流形，$G$ 是在 $M$ 上光滑可迁作用的 Lie 群。取定 $p\in M$，令 $G_p$ 为迷向群。则由 $F(gG_p)=g\cdot p$ 定义的映射 $F\colon G/G_p\to M$ 是 $G$-等变微分同胚。
\end{theorem}

\begin{proof}
$F$ 光滑，因为它由轨道映射 $\sigma_p\colon G\to M$ 取商得到（定理 4.30）。为证 $F$ 双射：任取 $q\in M$，由可迁性存在 $g\in G$ 使 $gp=q$，故 $F(gG_p)=q$。反之，若 $F(g_1H)=F(g_2H)$，则 $g_1p=g_2p$ 蕴含 $g_1^{-1}g_2 p=p$，故 $g_1^{-1}g_2\in G_p$，即 $g_1G_p=g_2G_p$。于是 $F$ 是等变光滑双射，由等变秩定理它是微分同胚。
\end{proof}

该定理表明齐性空间的研究可化归为理解 Lie 群对闭子群的商这一代数问题。把上述刻画定理应用于早先讨论的可迁群作用例子，可见一些熟悉空间微分同胚于 Lie 群对闭子群的商。

\begin{example}{[齐性空间再探]}\label{example:ch21-齐性空间再探}
(a) 再考虑 $O(n)$ 在 $\mathbb{S}^{n-1}$ 上的自然作用。若取北极 $N=(0,\dots,0,1)$ 为基点，易验迷向群为 $O(n-1)$（视为固定最后一个坐标的 $\mathbb{R}^n$ 正交变换）。故 $\mathbb{S}^{n-1}\cong O(n)/O(n-1)$。
(b) 对 $\operatorname{SO}(n)$ 在 $\mathbb{S}^{n-1}$ 上的作用，迷向群为 $\operatorname{SO}(n-1)$，故 $n\geqslant2$ 时 $\mathbb{S}^{n-1}\cong\operatorname{SO}(n)/\operatorname{SO}(n-1)$。
(c) 因 $E(n)$ 光滑可迁作用于 $\mathbb{R}^n$，原点迷向群为子群 $O(n)\subseteq E(n)$，故 $\mathbb{R}^n\cong E(n)/O(n)$。
(d) 再考虑 $\operatorname{SL}(2,\mathbb{R})$ 经 Möbius 变换在上半平面上的可迁作用。直接计算表明点 $\mathrm{i}\in U$ 的迷向群由形如 $\begin{pmatrix}\cos\theta&-\sin\theta\\ \sin\theta&\cos\theta\end{pmatrix}$ 的矩阵组成，即 $\operatorname{SO}(2)$。
\end{example}

由于齐性空间结构如此丰富，自然会问是否每个光滑流形都可实现为 Lie 群对闭子群的商，等价地，是否每个光滑流形都容许可迁 Lie 群作用。答案断然是否定的，因为存在很强的拓扑限制。首先，非连通流形的一个简单必要条件是其所有连通分支彼此微分同胚（习题 21-11）。更重要的是，G.~D.~Mostow 2005 年的一个引人注目的定理表明紧齐性空间必有非负 Euler 示性数（Euler 示性数定义见习题 18-9）。

\begin{theorem}{[带可迁群作用的集合]}\label{theorem:ch21-可迁作用集合}
设 $X$ 是集合，给定 Lie 群 $G$ 在 $X$ 上的可迁作用，且对某点 $p\in X$，迷向群 $G_p$ 在 $G$ 中闭。则 $X$ 有唯一光滑流形结构使该作用光滑。在此结构下 $\dim X=\dim G-\dim G_p$。
\end{theorem}

\begin{proof}
定理\cref{theorem:ch21-齐性空间构造} 表明 $G/G_p$ 是维数 $\dim G-\dim G_p$ 的光滑流形。由 $F(gG_p)=g\cdot p$ 定义的 $F\colon G/G_p\to X$，用与定理\cref{theorem:ch21-齐性空间刻画} 完全相同的论证是等变双射（该部分证明未用到 $M$ 是流形或作用光滑的假设）。令 $F$ 为微分同胚，在 $X$ 上定义拓扑与光滑结构，则 $G$ 在 $X$ 上的作用光滑，因为它可写成 $(g,x)\mapsto F(g\cdot F^{-1}(x))$。若 $\widetilde X$ 是使作用光滑的任一光滑流形结构，则刻画定理表明 $F$ 也是从 $G/G_p$ 到 $\widetilde X$ 的等变微分同胚，故 $\widetilde X$ 的拓扑与光滑结构等于上述构造。
\end{proof}

\begin{example}{[Grassmann 流形]}\label{example:ch21-Grassmann}
令 $G_k(\mathbb{R}^n)$ 表示 $\mathbb{R}^n$ 的 $k$ 维子空间的 Grassmann 流形如例 1.36。一般线性群 $\operatorname{GL}(n,\mathbb{R})$ 在 $G_k(\mathbb{R}^n)$ 上可迁：给定两子空间 $A_1$、$A_2$，为二者选基并扩充为 $\mathbb{R}^n$ 的基，把前者基映到后者基的线性变换即把 $A_1$ 映到 $A_2$。迷向群闭，故 $G_k(\mathbb{R}^n)$ 有唯一光滑流形结构使其成为齐性 $\operatorname{GL}(n,\mathbb{R})$-空间。
\end{example}

\begin{example}{[旗流形]}\label{example:ch21-旗流形}
设 $V$ 是 $n>1$ 维实向量空间，$K=(k_1,\dots,k_m)$ 为满足 $0<k_1<\cdots<k_m<n$ 的整数有限序列。$V$ 中类型 $K$ 的\emph{旗}是线性子空间的嵌套序列 $S_1\subset S_2\subset\cdots\subset S_m\subset V$，$\dim S_i=k_i$。$V$ 中所有类型 $K$ 的旗的集合记为 $F_K(V)$（例如 $K=(k)$ 时 $F_K(V)=G_k(V)$）。易证 $\operatorname{GL}(V)$ 在 $F_K(V)$ 上可迁，且迷向群为闭子群（习题 21-16），故 $F_K(V)$ 有唯一光滑流形结构使其成为齐性 $\operatorname{GL}(V)$-空间。在此结构下 $F_K(V)$ 称为\emph{旗流形}。
\end{example}
\end{document}
